\section{Introduction}
Multi-view LiDAR point cloud registration is a fundamental problem in robotics and computer vision, with applications in autonomous driving, SLAM, and 3D reconstruction. Traditional methods for LiDAR point cloud registration typically rely on feature-based or direct methods, which can be sensitive to noise and outliers.

Recently, Neural Radiance Fields (NeRFs)~\cite{mildenhall2020nerfOG} have gained significant interest for their powerful capabilities in scene reconstruction and novel view synthesis. A major advancement in RGB image NeRF (RGB NeRF) research is NeRF Bundle Adjustment (NeRF-BA)~\cite{lin2021barf,wang2022nerfmm,chen2023l2gnerf}, which leverages the differentiability of camera projection to jointly optimize the scene neural representation and camera poses. This has demonstrated remarkable success in recovering accurate camera poses from noisy pose initialization. Nonetheless, few studies have explicitly focused on the problem of LiDAR NeRF-BA, and the fundamental differences in the designing of RGB NeRFs and LiDAR NeRFs.  

Contrary to RGB images that contains rich texture for camera pose optimization in RGB NeRF-BA, LiDAR point clouds post the following challenges for LiDAR NeRF-BA: (1) LiDAR point clouds inherently lack textures. A direct implementation of RGB NeRF-BA on LiDAR point clouds would result in sub-optimal pose optimization, as it overlooks the strong geometric cue, i.e., accurate ranges, from LiDAR rays. (2) LiDAR point clouds are spatially sparse. Directly using the volume sampling strategy in RGB NeRFs would result in inefficient volume sampling and hinder pose optimization, which will be discussed in Section~\ref{sec:volume_sampling}.

To address these challenges, we investigate and show that LiDAR NeRF-BA performance is closely coupled with the density of volume sampling, and propose a novel volume sampling method tailored for LiDAR ray models to enhance LiDAR NeRF-BA. Given an initially noisy estimation of LiDAR poses in a sequence, we propose a novel LiDAR NeRF-BA algorithm, termed NeLD-BA, that jointly reconstruct scenes and refine LiDAR poses beyond the capabilities of existing LiDAR NeRF methods on novel view synthesis and odometry (e.g., Figure~\ref{fig:intro_example}). 

The main contributions of this work include: 1) We provide new insights that pose accuracy in LiDAR NeRF-BA is explicitly reinforced by volume sampling positions; 2) Novel, simple, yet efficient strategies are proposed for volume sampling and positional encoding tailored for LiDAR NeRF-BA with a LiDAR-specific loss function; 3) Extensive experiments are conducted against state-of-the-art algorithms, showing superior mapping and localization performance.
% Furthermore, we tackle the challenge of gradient instability caused by positional encoding~\cite{lin2021barf} by introducing surrogate gradients. Finally, inspired by DS-NeRF~\cite{deng2024dsnerf}, we design a loss function for LiDAR point clouds that enhances reconstruction quality and pose accuracy simultaneously.

\begin{figure}
  \centering
  \begin{subfigure}[t]{0.49\linewidth}
    \centering
    \includegraphics[width=\linewidth]{images/thumbnail/hba.png}
  \end{subfigure}\hfill
  \begin{subfigure}[t]{0.49\linewidth}
    \centering
    \includegraphics[width=\linewidth]{images/thumbnail/neld_ba.png}
  \end{subfigure}
  \caption{3D mapping results with HBA~\cite{liu2022HBA} (left) and the proposed NeLD-BA method (right) on the quad\_hard sequence from the Newer College dataset~\cite{Ramezani2020newercollege}. NeLD-BA produces a higher quality map.}
  \label{fig:intro_example}
  \vspace{-10pt}
\end{figure}

Finally, we will open-source our code for the community upon the acceptance of the paper. 
